% =============================================================================
% Project : THE_PHILOSOPHY_OF_VIABILITY
% Book    : Book I — Ontology of Viability
% Book ID : VS.B01
%
% File    : 4_ASSETS/diagrams/
%           VS.B01.P01.C03.S04.SS06.P004.DG01.identity_plurality_across_ontological_levels.tex
%
% Asset   : DG01 — Identity Plurality Across Ontological Levels
%
% Purpose :
%   To display distinct levels at which unity and plurality
%   may be asserted, without implying a necessary hierarchy
%   or numerical correspondence among them.
%
% Requires (preamble):
%   \usepackage{tikz}
%   \usetikzlibrary{positioning, fit}
% =============================================================================

\begin{figure}[htbp]
    \centering

    \begin{tikzpicture}[
        font=\small,
        node distance=0.4cm,
        % --- base box ---
        box/.style={
            draw, rounded corners=3pt, align=center,
            inner sep=6pt, line width=0.7pt
        },
        % --- node roles ---
        levelbox/.style={box, text width=2.8cm, minimum height=1.4cm,
                         fill=blue!8, draw=blue!55!black},
        identitybox/.style={box, text width=6.4cm, minimum height=1.3cm,
                            fill=blue!16, draw=blue!55!black},
        % --- frames and labels ---
        vsgroup/.style={draw=black!45, dashed, rounded corners=5pt,
                        inner sep=8pt, line width=0.5pt},
        grouplabel/.style={font=\footnotesize\itshape, text=black!75,
                           align=left},
        elabel/.style={font=\footnotesize\itshape, text=black!75,
                       align=left, text width=4.6cm},
        % --- undirected link (no arrowheads: no priority, no entailment) ---
        vslink/.style={draw=black!55, dashed, line width=0.7pt}
    ]

    % ---------------------------------------------------------------
    % Row 1: three levels
    % ---------------------------------------------------------------
    \node[levelbox] (manifestation) {
        \textbf{Manifestation}\\[3pt]
        One or many forms
    };
    \node[levelbox, right=of manifestation] (components) {
        \textbf{Components}\\[3pt]
        One or many elements
    };
    \node[levelbox, right=of components] (relations) {
        \textbf{Relations}\\[3pt]
        One or many connections
    };

    % ---------------------------------------------------------------
    % Row 2: two levels, centred under the row above
    % ---------------------------------------------------------------
    \node[levelbox, below=of components, xshift=-1.8cm] (source) {
        \textbf{Source}\\[3pt]
        One or many grounds
    };
    \node[levelbox, below=of components, xshift=1.8cm] (aggregate) {
        \textbf{Aggregate}\\[3pt]
        Unity by a grouping criterion
    };

    % ---------------------------------------------------------------
    % Frame: five levels of equal rank
    % ---------------------------------------------------------------
    \node[vsgroup, fit=(manifestation)(components)(relations)
                       (source)(aggregate)] (levels) {};
    \node[grouplabel, anchor=south west] at (levels.north west) {
        Levels at which unity and plurality may be asserted
    };

    % ---------------------------------------------------------------
    % Identity (outside the frame)
    % ---------------------------------------------------------------
    \node[identitybox, below=1.1cm of levels] (identity) {
        \textbf{Identity}\\[3pt]
        Unity or plurality requiring identity-level grounds
    };

    \draw[vslink] (levels.south) -- (identity.north)
        node[elabel, midway, right=5pt] {
            possible relevance for inquiry
        };

    % ---------------------------------------------------------------
    % Analytical constraint
    % ---------------------------------------------------------------
    \node[below=0.7cm of identity, text width=10.5cm, align=center,
          font=\footnotesize\itshape] {
        \textbf{Analytical constraint:}
        A unity or plurality claim at one level does not determine
        the identity configuration at another level without
        additional ontological grounds. The connection indicates
        possible relevance for inquiry, not entailment, causation,
        or a fixed hierarchy.
    };

    \end{tikzpicture}

    \caption{Distinct ontological levels relevant to the analysis of
    identity plurality}
    \label{fig:identity-plurality-across-ontological-levels}
\end{figure}

% =============================================================================
% END OF ASSET DG01
% =============================================================================